% Overview of RTGS Framework — compact (~1/4 page)
% IEEE Transactions style | third-person passive

\subsection{Overview of the RTGS Framework}
\label{subsec:overview}

Conventional offloading frameworks treat the microservice DAG
$\mathcal{G}_t = (\mathcal{V}_t, \mathcal{E}_t)$ as a fixed input
and optimise only the task-to-UAV assignment
$\pi^*: \mathcal{V}_t \to \mathcal{V}_r$.
This decoupled strategy is inherently limited in heterogeneous UAV
swarms: as UAVs move, link bandwidths $B_{kl}(n)$ fluctuate and
battery depletion shrinks $\tilde{\phi}_k(n)$, so a topology that
yields a feasible schedule at epoch $n$ may violate the deadline at
epoch $n{+}1$—even under an optimal fixed-graph assignment policy.

RTGS addresses this limitation by promoting the graph topology itself
to a control variable, seeking the joint solution
\begin{equation}
  (\mathcal{G}_t^{\ast},\,\pi^{\ast}) =
  \arg\min_{\mathcal{G}_t' \in \Omega,\;\pi}\;
  T_{\mathrm{span}}\!\bigl(\mathcal{G}_t',\,\pi,\,\mathcal{G}_r(n)\bigr),
  \label{eq:joint_opt}
\end{equation}
where $\Omega$ is the set of DAGs reachable from $\mathcal{G}_t$ via
the \textsc{Split} and \textsc{Merge} operators defined in
Section~\ref{subsubsec:operators}.
As illustrated in Fig.~\ref{fig:framework}, RTGS realises
\eqref{eq:joint_opt} through three tightly coupled components:
(i)~a \textbf{dual-graph GNN} with cross-attention encodes the
joint state of $\mathcal{G}_t$ and $\mathcal{G}_r(n)$ into
resource-aware per-node embeddings;
(ii)~a \textbf{node-level PPO agent} selects a graph-shaping action
from a 211-dimensional space ($1$ noop, $20$ split, $190$ merge)
conditioned on those embeddings;
and (iii)~a \textbf{deterministic greedy scheduler} maps the
reshaped graph $\mathcal{G}_t'$ to a concrete UAV assignment,
whose outcome closes the feedback loop as the reward signal.
Because the scheduling layer is non-differentiable, RL is the
natural vehicle for optimising the topology decisions end-to-end.
